%% ---------------------------------------------------------------------
\section{La fórmula de Ackermann}
\label{sec:re-ackermann}
%% ---------------------------------------------------------------------

Comparar coeficientes funciona bien con $n=2$ o $3$, pero con más estados el álgebra se vuelve pesada. Para sistemas de una entrada hay una fórmula cerrada.

\begin{historia}[Jürgen Ackermann]
La fórmula lleva el nombre del ingeniero alemán Jürgen Ackermann, que la publicó en 1972 en un artículo sobre el diseño de sistemas de control lineales en el espacio de estado~\citep{Ack72}. Ackermann trabajó luego en el Centro Aeroespacial Alemán (DLR) en control robusto, con aplicaciones a la dirección de automóviles.
\end{historia}

\begin{definicion}[label=def:re-ackermann]{Fórmula de Ackermann}
Si el sistema de una entrada es controlable y el polinomio característico deseado es $q(\lambda)=\lambda^n+\alpha_{n-1}\lambda^{n-1}+\cdots+\alpha_1\lambda+\alpha_0$, entonces
\begin{equation}
  \Km=\bigl[\,0\;\;0\;\;\cdots\;\;0\;\;1\,\bigr]\,\Ctrb^{-1}\,q(\Am),
  \qquad
  q(\Am)=\Am^n+\alpha_{n-1}\Am^{n-1}+\cdots+\alpha_1\Am+\alpha_0\Id.
  \label{eq:re-ackermann}
\end{equation}
\end{definicion}

Observe que $q(\Am)$ es el polinomio deseado evaluado en la \emph{matriz} $\Am$ (con $\Id$ en lugar del término independiente), y que la fórmula exige $\Ctrb^{-1}$: solo funciona si el sistema es controlable.

\begin{ejemplo}[label=ej:re-ackermann]{El doble integrador}
Para $Y(s)/U(s)=1/s^2$, halle $\Km$ para tener los polos en $-1\pm j$.
\solucion En variables de fase, $\Am=\begin{bmatrix}0&1\\0&0\end{bmatrix}$, $\Bm=\begin{bmatrix}0\\1\end{bmatrix}$. El polinomio deseado es $q(\lambda)=\lambda^2+2\lambda+2$. Entonces
\[
  \Ctrb=[\Bm\;\;\Am\Bm]=\begin{bmatrix}0&1\\1&0\end{bmatrix}=\Ctrb^{-1},\qquad
  q(\Am)=\Am^2+2\Am+2\Id=\begin{bmatrix}2&2\\0&2\end{bmatrix},
\]
\[
  \Km=[0\;\;1]\begin{bmatrix}0&1\\1&0\end{bmatrix}\begin{bmatrix}2&2\\0&2\end{bmatrix}=[1\;\;0]\begin{bmatrix}2&2\\0&2\end{bmatrix}=[2\;\;2].
\]
Comprobación: $\Am-\Bm\Km=\begin{bmatrix}0&1\\-2&-2\end{bmatrix}$, cuyo polinomio es $\lambda^2+2\lambda+2$.
\end{ejemplo}

\codigorepo{cap09_retroalimentacion_estado/C9_Ej4_Ackermann.m}{ackermann}{ejemplo \ref{ej:re-ackermann}: la fórmula paso a paso, \cmd{acker} y \cmd{place}}

\begin{cuidado}[\texttt{acker} o \texttt{place}]
\cmd{acker} implementa \eqref{eq:re-ackermann}: es exacta en papel pero numéricamente frágil para $n$ grande (invierte $\Ctrb$, que puede estar mal condicionada). \cmd{place} usa un algoritmo más robusto, sirve para varias entradas, pero no acepta polos repetidos más veces que el número de entradas. Para los ejemplos del curso dan lo mismo.
\end{cuidado}

\begin{receta}[Ubicación de polos por retroalimentación de estado]
\begin{enumerate}
  \item Escribir el modelo de estado y verificar que $\operatorname{rango}\Ctrb=n$.
  \item Traducir las especificaciones en polos deseados (\cref{sec:re-especificaciones}) y armar $q(\lambda)$.
  \item Calcular $\Km$: comparando los coeficientes de $\det(\lambda\Id-\Am+\Bm\Km)$ con $q(\lambda)$, o con \eqref{eq:re-ackermann}.
  \item Verificar \cmd{eig(A-B*K)}, simular, y mirar el esfuerzo de control $u(t)$.
\end{enumerate}
\end{receta}

\begin{pausa}
Use la fórmula de Ackermann para $1/(s(s+3))$ en variables de fase, con polos en $-4\pm j4$.
\respuesta{$\Am=\begin{bmatrix}0&1\\0&-3\end{bmatrix}$, $\Ctrb=\begin{bmatrix}0&1\\1&-3\end{bmatrix}$, $\Ctrb^{-1}=\begin{bmatrix}3&1\\1&0\end{bmatrix}$, $q(\Am)=\Am^2+8\Am+32\Id=\begin{bmatrix}32&5\\0&17\end{bmatrix}$. Entonces $\Km=[1\;\;0]\,q(\Am)=[32\;\;5]$. Comprobación: $\lambda^2+(3+K_2)\lambda+K_1=\lambda^2+8\lambda+32$.}
\end{pausa}
